\section{Introduction}
\label{sec:intro}
Cooperative marine monitoring assigns spatially dispersed tasks with varied payload requirements to heterogeneous platforms, as in air--sea environmental monitoring.\cite{bi2024} A feasible allocation must match tasks to capabilities, respect service-start time windows, and retain sufficient transfer energy for subsequent visits and return to base. These constraints couple task-to-platform assignment with visit order, requiring allocation and route planning to be considered jointly.\cite{liu2018}

Constrained task allocation requires feasible assignments and coordinated execution.\cite{nunes2017} Auctions coordinate assignments through bids and task bundles,\cite{choi2009} while evolutionary search addresses capability and timing restrictions for heterogeneous surface vehicles.\cite{wang2024} Adaptive large neighborhood search (ALNS) improves routes by removing and reinserting tasks, adapting operator sampling according to past rewards.\cite{ropke2006} These rewards describe which changes have been productive, whereas current task insertability identifies requests that remain difficult to allocate. The two sources of information are complementary and can jointly direct route reconstruction.

Learning-based search relates the evolving solution to neighborhood control. Graph networks have selected removals while an optimization solver repairs routes,\cite{zhou2023} and reinforcement learning has adapted ALNS operators, removal strength, and acceptance settings.\cite{reijnen2024} Building on this state-driven approach, a large language model (LLM) can interpret allocation observations alongside previous actions and outcomes, then propose controls that a route solver executes while enforcing feasibility.

LLMs have been used to generate candidate solutions from previous evaluations,\cite{opro2024} improve executable heuristics through reflective feedback,\cite{reevo2024} and develop scheduling rules that adjust neighborhood weights during search.\cite{ns4s2025} In marine task allocation, the changing relationship between unassigned tasks and platform routes calls for translating allocation difficulties into a search focus and executable neighborhood operations.

We propose IntALNS, Interpretable Agent-Controlled Adaptive Large Neighborhood Search, to generate state-dependent neighborhood controls and explicit rationales from allocation observations, execution history, and strategy examples. The LLM agent interprets the evolving allocation state and directs search, while ALNS reconstructs feasible routes and returns execution feedback. Experiments evaluate allocation quality and search progress on synthetic marine monitoring instances; a recorded case examines the link between the observed state, agent rationale, executed control, and subsequent outcome.
